\chapter{研究结论与展望}
\label{chap:conclusion}

\section{主要结论}

基于前文的测度与分解，本文得到以下主要结论：

\begin{enumerate}[label=（\arabic*）]
    \item 样本期内数字经济发展水平总体呈上升趋势，但地区间绝对差距仍在扩大；
    \item 区域差异主要来源于组间差距，东部与中西部之间的差距是总差异的主要构成；
    \item 数字基础设施与数字产业化维度对综合得分的贡献相对更大。
\end{enumerate}

\section{政策建议}

针对上述结论，本文提出三点建议：加大中西部数字基础设施投入以补齐短板；推动区域间产业协同与数据要素流动；完善数字治理的制度供给，营造良好的数字经济发展环境。

\section{不足与展望}

受演示数据所限，本文结论仅用于说明方法流程。未来可在真实数据基础上，引入空间计量模型刻画区域间溢出效应，并延长样本期以考察数字经济发展的动态演变趋势。
